\documentclass[conference]{IEEEtran}
\IEEEoverridecommandlockouts

\usepackage{cite}
\usepackage{amsmath,amssymb,amsfonts}
\usepackage{booktabs}
\usepackage{graphicx}
\usepackage{textcomp}
\usepackage{xcolor}
\usepackage{url}
\usepackage{array}
\newcolumntype{L}[1]{>{\raggedright\arraybackslash}p{#1}}
\newcommand{\figplaceholder}[2]{%
\fbox{\begin{minipage}{0.92\linewidth}
\centering
\vspace{0.5em}
\textbf{Figure placeholder: #1}\\[0.35em]
\small #2
\vspace{0.5em}
\end{minipage}}}

\def\BibTeX{{\rm B\kern-.05em{\sc i\kern-.025em b}\kern-.08em
    T\kern-.1667em\lower.7ex\hbox{E}\kern-.125emX}}

\begin{document}

\title{From Quality Filtering to Curation Decisions: A Stage-Boundary Framework for Language-Model Training Data}

\author{\IEEEauthorblockN{Keonsoo Lee}
\IEEEauthorblockA{\textit{Department of Computer Science} \\
\textit{University of Nevada, Las Vegas}\\
Las Vegas, USA \\
leek133@unlv.nevada.edu}
\and
\IEEEauthorblockN{Shaikh Arifuzzaman}
\IEEEauthorblockA{\textit{Department of Computer Science} \\
\textit{University of Nevada, Las Vegas}\\
Las Vegas, USA \\
shaikh.arifuzzaman@unlv.edu}
}

\maketitle

\begin{abstract}
Foundation-model training pipelines increasingly depend on large,
heterogeneous, and repeatedly refreshed candidate corpora.  The operational
step between data collection and model training is often described as generic
filtering or quality scoring, although real training-data decisions combine
structural validity, risk quarantine, redundancy control, coverage
preservation, budget allocation, and downstream validation.  This paper argues
that training-data preparation should be treated as a curation decision
problem rather than a single quality-filtering problem.  The framework
separates operational responsibilities into Core-Metric-Policy layers and
enforces a staged contract:
Stage 0 quarantines risk, Stage A applies chunk-level hard gates, Stage B
performs optional budgeted chunk selection, and Stage C validates the resulting
subset through held-out training evidence.  A central design constraint is that
Utility is measured only in Stage C and is never used as a Stage-B selector
objective.  The novelty is not another quality score; it is an auditable
separation between record disposition, budget allocation, and downstream
validation.  The current evidence contains one positive code-domain validation
and one negative math-domain validation.  In code, a natural-budget curated
arm reduces packed training tokens by 60.8\% while improving held-out NLL and
EvalPlus macro pass rate over raw fine-tuning.  In math, the first selector
fails, and the repaired selector improves over that failure but still does not
beat raw training.  A frozen 48-task LiveCodeBench pilot gives identical
18.75\% pass@1 to base, raw, and curated arms, so independent benchmark
transfer is not yet demonstrated.  The result is not a universal data-quality detector or a
production-ready safety filter; it is evidence that curation policies should
expose pass, fail, and abstain outcomes under frozen validation rather than
collapse all data preparation into an opaque quality score.
\end{abstract}

\begin{IEEEkeywords}
big data, data curation, foundation models, language models, pretraining data,
data-centric AI, data selection, deduplication, reproducibility
\end{IEEEkeywords}

\section{Introduction}
Training data is a primary control surface for language models.  As model
training and updating cycles continue, practitioners may collect new candidate
data between training runs, from web corpora, code repositories, documentation,
domain-specific sources, or generated and human-authored mixtures.  The central
question is not whether every collected item is intrinsically ``high quality.''
Rather, the operational question is whether the collected corpus can be turned
into a training-ready pool, or an explicitly budgeted training subset, under a
reproducible and auditable policy.

This distinction matters because raw candidate corpora can fail in different
ways.  Some records are structurally unusable, contaminated, unsafe, malformed,
or outside the intended rights boundary.  Some records are individually usable
but redundant enough to waste a constrained training budget.  Some records are
valuable precisely because they preserve rare styles, sources, domains, or
task forms.  A single scalar quality score cannot represent all of these
decisions without hiding the policy assumptions behind the score.

We therefore frame data preparation as a \emph{curation stage} between
collection and training.  The curation stage does not claim to solve universal
data quality.  It maps candidate data into explicit dispositions: quarantine,
Stage-A pass, retained curated pool, optional budget-selected subset, or
abstention when the evidence is insufficient.  This paper focuses on that
operational layer.  This changes the research question from ``Can we score data
quality?'' to ``Can we make training-data decisions reproducible, stage-local,
and falsifiable?''

The main contributions are:
\begin{itemize}
    \item \textbf{A decision framework, not a quality score.}  We define a
    curation contract that separates record disposition, optional budget
    allocation, and downstream validation.
    \item \textbf{A Core-Metric-Policy method.}  We map each Core
    responsibility to metric surfaces, stage ownership, and forbidden claims.
    \item \textbf{A Utility isolation rule.}  We enforce that downstream
    Utility is measured only in Stage C and is never available to the Stage-B
    selector objective.
    \item \textbf{An evidence package with positive and abstain outcomes.}  We
    report historical positive code-domain evidence awaiting a current-framework
    rerun, alongside a math-domain abstain under the same frozen validation
    discipline.
\end{itemize}

This framing gives reviewers a falsifiable target.  The claim is not that the
framework can certify any dataset as high quality.  The claim is that an
LM-training corpus can be processed through explicit curation stages whose
outputs, failures, and unsupported claims are machine-checkable.

\begin{figure*}[t]
\centering
\figplaceholder{Curation-stage pipeline}{Replace with a full-width pipeline from raw candidate corpus to Stage 0, Stage A, full curated pool, optional Stage B, and Stage C.  Mark Stage C as validation-only and show the blocked Utility feedback path into Stage B.}
\caption{Curation-stage view of language-model training data management.  Stage B is optional and budget-driven; Utility is observed only in Stage C and cannot flow back into the Stage-B selector objective.}
\label{fig:pipeline}
\end{figure*}

\section{Motivation and Scope}
Large-scale language-model corpora commonly use filtering, deduplication,
contamination checks, source balancing, and quality-related heuristics.  Prior
work on web-scale corpus construction and dataset filtering motivates these
operations, while deduplication studies show that repeated text can affect both
evaluation and training behavior~\cite{ccnet,dedup,refinedweb,dolma}.  However,
many practical systems combine several different objectives under a broad
quality-filtering label.

The big-data aspect is the operational state of the corpus, not token count
alone.  Candidate corpora are large, heterogeneous, refreshed over time, and
useful only after admissibility, redundancy, coverage, budget, and downstream
evidence are recorded.  We therefore treat curation as an auditable
data-management stage before training, rather than as a hidden preprocessing
detail inside a model-training recipe.

Our scope is narrower and more explicit.  We do not attempt to define an
intrinsic semantic quality oracle.  We also do not claim production-grade
license, PII, secret, poisoning, or benchmark-contamination detection.  Instead,
we study whether a collected corpus can be passed through a documented curation
stage whose gates, selector, diagnostics, and validation evidence are
reproducible.

The current evidence supports a bounded research claim: the proposed
curation-stage decision structure can be specified, implemented, audited, and
falsified under frozen downstream validation.  It does not support a production
deployment claim or a universal all-domain improvement claim.

\section{Related Work}
\subsection{Large-Scale Corpus Construction}
Large language-model corpora are usually built through layered data-processing
pipelines rather than through a single quality decision.  C4 popularized a
large cleaned Common Crawl corpus for transfer learning experiments~\cite{c4}.
CCNet used language identification, deduplication, and filtering to extract
monolingual web corpora~\cite{ccnet}.  The Pile emphasized mixture diversity
across multiple high-quality sources~\cite{pile}.  RefinedWeb showed that
filtered and deduplicated web data can be competitive at large scale~\cite{refinedweb}.
Dolma documented an open three-trillion-token corpus and its curation
practices~\cite{dolma}, while FineWeb-Edu further highlighted how filtering
choices can change downstream behavior~\cite{finewebedu}.

These projects motivate the need for data curation, but they do not by
themselves define a reusable policy boundary between hard gates, budgeted
selection, and downstream validation.  Our work contributes that policy
separation and the associated reproducibility package.

\subsection{Deduplication, Diversity, and Data-Centric Evaluation}
Deduplication is a necessary but incomplete curation operation.  Lee et
al. showed that duplicated training data affects memorization and evaluation
overlap~\cite{dedup}.  However, for code and technical text, repeated APIs,
tests, and idioms can also be useful recurrence.  The framework therefore
separates hard duplicate gates from softer redundancy-risk penalties.

Data-centric benchmarks such as DataComp make the dataset construction process
itself an object of evaluation~\cite{datacomp}.  This paper follows the same
spirit for language-model training data: the curation stage is evaluated as a
reproducible data-management process, not merely as a hidden preprocessing
detail before model training.

\section{Problem Setting}
Let $D=\{x_i\}_{i=1}^{n}$ be a candidate corpus collected before a language-model
training run.  The goal is not to assign a universal quality score to every
$x_i$.  The goal is to produce a reproducible disposition for each record and,
when a token budget $B$ is binding, a training subset $S_B \subseteq D$ whose
construction is auditable.

The framework must support five outcomes.  First, a record can be quarantined
because it violates a risk or admissibility boundary.  Second, it can fail
Stage A and be excluded from the curated pool.  Third, it can pass Stage A and
remain in the full curated pool.  Fourth, it can be selected by Stage B under a
budgeted training subset.  Fifth, the whole subset can be held, rejected, or
accepted after Stage-C validation.  These outcomes are different policy states:
``budget not selected'' is not the same as ``bad data,'' and a retain-all
outcome is valid when the corpus already fits the training budget.

This formulation yields three requirements:
\begin{enumerate}
    \item decisions must be stage-local, so that chunk-level gates do not
    accidentally consume subset-level outcomes;
    \item metrics must be attached to explicit policy roles rather than
    interpreted as free-standing quality claims;
    \item each claim must be traceable to frozen artifacts, commands, configs,
    and validation reports.
\end{enumerate}

\section{Framework}
\subsection{Core-Metric-Policy Separation}
The framework distinguishes three layers.  A \emph{Core} is an operational
curation responsibility.  A \emph{Metric} is an observable proxy or measurement
used to support that responsibility.  A \emph{Policy} is the stage-specific rule
that consumes the metric.  This separation prevents a metric from silently
expanding into a claim it cannot support.

The current Core surfaces are:
\begin{itemize}
    \item \textbf{Validity:} structural usability and basic training-input
    admissibility;
    \item \textbf{Selection Value Evidence:} observable pre-outcome evidence
    used for optional budget allocation, not intrinsic quality judgment;
    \item \textbf{Redundancy:} duplicate and saturation control, while
    distinguishing harmful duplication from useful recurrence;
    \item \textbf{Coverage:} observable retention of source, style, path,
    content type, and cluster structure;
    \item \textbf{Utility:} downstream training effect, measured only after a
    subset is constructed.
\end{itemize}

The term Quality is retained only as a legacy implementation label where
necessary.  The paper-level construct is Selection Value Evidence, because the
framework cannot prove that a chunk has intrinsic semantic quality from
automatic proxy features alone.

\begin{table*}[t]
\caption{Core-Metric-Policy Mapping}
\label{tab:corepolicy}
\centering
\small
\begin{tabular}{L{0.13\textwidth}L{0.30\textwidth}L{0.28\textwidth}L{0.19\textwidth}}
\toprule
Core & Operational responsibility & Example metric surface & Policy role \\
\midrule
Validity & Structural usability and admissibility as pretraining input & structural validity gate, malformed-text checks & Stage-A hard gate \\
Selection Value Evidence & Observable pre-outcome evidence for budget allocation; not intrinsic quality & code usefulness proxy, source/path/style support & Stage-B ranking signal when budget binds \\
Redundancy & Duplicate and saturation control while preserving useful recurrence & canonical-exact duplicate indicator, fuzzy near-duplicate risk & Exact-only Stage-A gate and reversible Stage-B penalty \\
Coverage & Retention of observable source, style, path, content-type, and cluster structure & coverage retention and tail/source diagnostics & Stage-C subset validator and diagnostic \\
Utility & Downstream training effect under frozen comparison arms & held-out NLL and guardrails & Stage-C validation only \\
\bottomrule
\end{tabular}
\end{table*}

\begin{figure}[t]
\centering
\figplaceholder{Core-Metric-Policy map}{Replace with a compact matrix linking Validity, Selection Value Evidence, Redundancy, Coverage, and Utility to metric surfaces, stage ownership, and claim boundaries. Highlight Utility as Stage C only.}
\caption{The framework treats metrics as policy-bound evidence, not as free-standing quality claims.}
\label{fig:corepolicy}
\end{figure}

\subsection{Stages}
The curation pipeline is organized into four stages.

\textbf{Stage 0: risk quarantine.}  Stage 0 removes or quarantines candidates
that violate project-defined safety or admissibility boundaries, including
detected secrets, PII-like records, benchmark-contamination indicators,
poisoning patterns, malformed text, or unsupported rights status.  In the
current work, this is a scoped project detector, not a production compliance
system.

\textbf{Stage A: chunk-level hard gate.}  Stage A applies structural validity
and hard duplicate gates.  A Stage-A failure means the chunk should not enter
the curated pool under the current contract.

\textbf{Stage B: chunk-level selection.}  Stage B ranks Stage-A-passing chunks
only when a binding training budget requires a subset.  It may use pre-outcome
selection evidence and redundancy risk, but it cannot use Stage-C Utility or
training outcomes.  If no budget requires reduction, retaining all Stage-A
passing chunks is valid.

\textbf{Stage C: subset-level validation.}  Stage C evaluates the constructed
subset.  Utility is measured here through downstream training evidence and
guardrails.  Stage-C outcomes may support, reject, or abstain from a paper
claim, but they are not fed back into Stage-B selection on frozen holdouts.

\begin{table*}[t]
\caption{Stage Responsibilities and Forbidden Uses}
\label{tab:stages}
\centering
\begin{tabular}{p{0.08\textwidth}p{0.41\textwidth}p{0.41\textwidth}}
\toprule
Stage & Allowed responsibility & Forbidden use \\
\midrule
0 & Quarantine scoped risk and admissibility hazards & Production legal or safety certification \\
A & Chunk-level hard gate for structural failures and hard duplicates & Utility-aware filtering \\
B & Optional budgeted selection among Stage-A-passing chunks & Consuming Stage-C outcomes \\
C & Subset validation through training evidence and guardrails & Retuning Stage B on frozen holdouts \\
\bottomrule
\end{tabular}
\end{table*}

\begin{figure}[t]
\centering
\figplaceholder{Utility leakage boundary}{Replace with a two-side diagram: allowed Stage-B selector inputs on the left, Stage-C outcomes on the right, and a blocked feedback arrow from Stage C back to Stage B.}
\caption{Utility is downstream validation evidence and is forbidden as a Stage-B selector objective.}
\label{fig:utilityboundary}
\end{figure}

\section{Implementation}
The implementation records each major decision as a machine-readable artifact.
Core behavior audits check that each Core surface remains within its allowed
interpretation.  Selector-leakage audits check that Utility-like fields do not
enter Stage-B selector objectives.  Comparison tables freeze training arms,
token budgets, and ablations before Stage-C interpretation.  A reproducibility
manifest records scripts, configs, outputs, and local hardware policy.

The artifact structure is grouped by scientific role: Core defense, Stage-C
validation, decision boundary, comparison tables, and reproducibility.  This
organization is deliberately different from a training log.  The goal is to
make each claim traceable to the stage where it is allowed, and to make failed
or incomplete evidence visible rather than treating it as a hidden engineering
detail.

\subsection{Audit Boundaries}
The audits are intentionally conservative.  A passing operational audit means
that the role declaration and stage boundary are consistent.  It does not mean
that every metric is a validated psychological, semantic, or legal construct.
Similarly, a Stage-0 detector benchmark shows scoped fixture behavior; it does
not certify a production detector.  The evidence artifacts therefore contain
both supported claims and forbidden claims.

\subsection{Selector Leakage Control}
The framework treats Utility as a post-selection outcome.  Stage-B inputs are
pre-outcome signals only.  Leakage control is therefore not a cosmetic
engineering check; it is required for the scientific claim.  If Stage-C NLL
results were allowed to tune Stage-B on the same frozen holdouts, the
comparison would collapse into outcome-driven optimization rather than a
curation-stage policy evaluation.

\section{Experimental Setup}
\subsection{Corpus Scenarios}
The validation setting contains one positive code-domain case and one negative
math-domain case.  Code is a useful stress case because code corpora mix
natural language, structured symbols, repeated APIs, tests, templates, project
metadata, and repository-specific conventions.  Math is a useful counter-case
because short answers, long derivations, worked examples, noisy extraction
artifacts, and rare topics can create different failure modes for the same
selection logic.  The experiment therefore evaluates not only whether curation
can improve a training arm, but also whether the framework exposes a failed
domain instead of hiding it.

\subsection{Training Arms}
The current natural-budget comparison reports base, raw, and curated arms.  The
base arm checks whether fine-tuning changes the model at all.  The raw arm
uses the collected corpus without Stage-B selection.  The curated arm uses the
framework output under the frozen policy.  For earlier equal-token development
runs, Stage-A random and known-reference arms provide additional diagnostic
baselines, but the current paper table emphasizes the natural-budget
comparison because it matches the practical question: can the framework train
less while preserving or improving downstream behavior?

The Stage-B ablation table also compares the full selector with quality-only,
redundancy-only, and no-coverage-support variants.  These ablations are
diagnostic; they are not Utility evidence by themselves.

\subsection{Evaluation}
The primary Stage-C measurement is held-out negative log-likelihood (NLL);
lower is better.  The code case also uses EvalPlus macro pass rate and a frozen
LiveCodeBench pilot as external capability checks.  The math case is currently
evaluated by held-out NLL only and is reported as failed.  All Stage-C outcomes
are validation evidence, not selector features.

The primary estimand is the mean NLL reduction
\begin{equation}
\Delta_{\mathrm{curated},b}=\frac{1}{K}\sum_{k=1}^{K}
\left(\mathrm{NLL}_{b,k}-\mathrm{NLL}_{\mathrm{curated},k}\right),
\end{equation}
where $b$ is a baseline arm and $K$ is the number of frozen seeds.  A positive
value means the curated arm has lower NLL than the baseline.  The natural-budget
table reports raw as the main practical baseline; equal-token development runs
retain Stage-A random as a diagnostic baseline for isolating the Stage-B effect.

\section{Results}
Table~\ref{tab:stagec} summarizes the current natural-budget Stage-C evidence.
Lower NLL is better.  Code is positive on five-seed held-out NLL and EvalPlus,
while an independent LiveCodeBench pilot is neutral.  Math v2 is the failed
selector, and math v3 is the repair-only abstain case.

\begin{table}[t]
\caption{Frozen Natural-Budget Stage-C Evidence}
\label{tab:stagec}
\centering
\scriptsize
\setlength{\tabcolsep}{2.5pt}
\begin{tabular}{llrccc}
\toprule
Domain & Arm & Tokens & NLL & Eval+ & Decision \\
\midrule
Code & Base & 0 & 1.23418 & 47.18 & Ref. \\
Code & Raw & 980992 & 1.20998 & 51.08 & Ref. \\
Code & Curated & 385024 & 1.20090 & 58.22 & Pass \\
Math & Base & 0 & 1.55945 & -- & Ref. \\
Math & Raw & 1120256 & 1.49565 & -- & Ref. \\
Math & Curated v2 & 626688 & 1.52706 & -- & Fail \\
Math & Curated v3 & 1026048 & 1.49899 & -- & Abstain \\
\bottomrule
\end{tabular}
\end{table}

In code, the curated arm uses 60.8\% fewer packed training tokens than raw,
improves held-out NLL by 0.00908, and improves same-protocol EvalPlus macro pass rate from
51.08\% to 58.22\%.  The frozen LiveCodeBench seed-101 pilot does not reproduce
an external capability gain: base, raw, and curated each pass 9 of 48 tasks.
Generated programs differ across arms, but the paired correctness vector is
identical.  In math, selector v2 uses 44.1\% fewer packed tokens but
worsens held-out NLL by 0.03141.  Selector v3 repairs most of that regression,
using 8.4\% fewer packed tokens than raw and reducing NLL to 1.49899, but it
still does not beat raw.  The paper therefore reports math as an abstain case,
not as a success and not as a tuning target.

\begin{figure}[t]
\centering
\figplaceholder{Natural-budget result summary}{Replace with a two-panel chart. Panel A: Code base/raw/curated NLL and EvalPlus. Panel B: Math base/raw/curated-v2/curated-v3 NLL. Use subdued colors and mark lower NLL as better.}
\caption{Code NLL and EvalPlus evidence is positive, the LiveCodeBench pilot is neutral, and Math remains an abstain case.}
\label{fig:results}
\end{figure}

Table~\ref{tab:ablation} summarizes the Stage-B ablation diagnostics.  The full
selector selects 423 chunks with a token proxy of 120235 and lower mean soft
redundancy risk than the quality-only variant.  These diagnostics support the
interpretation that the selector is not merely a quality-only ranking rule.
The quality-only variant selects slightly more chunks but carries higher
redundancy risk, while the redundancy-only variant aggressively reduces
redundancy at the cost of selecting far fewer chunks.  The full selector
therefore occupies the intended middle ground: it prefers concise and useful
chunks while retaining enough coverage support to avoid a pure deduplication
policy.  A separate Stage-B feature-shift diagnostic is consistent with this
mechanism: among 1913 Stage-A-passing records, the 1424 selected records have
higher structural richness and lower soft redundancy risk than the 489
budget-not-selected records.

\begin{table}[t]
\caption{Stage-B Ablation Diagnostics}
\label{tab:ablation}
\centering
\footnotesize
\setlength{\tabcolsep}{3pt}
\begin{tabular}{lccc}
\toprule
Arm & Chunks & Token proxy & Redundancy risk \\
\midrule
Full selector & 423 & 120235 & 0.19243 \\
Quality only & 446 & 120238 & 0.22203 \\
Redundancy only & 270 & 120238 & 0.04087 \\
No coverage support & 412 & 120237 & 0.18621 \\
\bottomrule
\end{tabular}
\end{table}

\begin{table}[t]
\caption{Claim Boundary from the Current Evidence Package}
\label{tab:claimboundary}
\centering
\footnotesize
\begin{tabular}{@{}p{0.46\columnwidth}p{0.48\columnwidth}@{}}
\toprule
Supported bounded claim & Unsupported stronger claim \\
\midrule
Historical Code Stage-C result; current rerun required & Universal domain improvement \\
Math abstain remains reported & Hiding failed domains \\
Utility is post-selection evidence & Utility-driven Stage-B selection \\
Core behavior is scoped and audited & Proof of intrinsic quality \\
Artifact set is reproducible & Production safety certification \\
\bottomrule
\end{tabular}
\end{table}

\begin{figure}[t]
\centering
\figplaceholder{Evidence and claim boundary}{Replace with a stacked evidence diagram: commands/configs, scoring artifacts, audits, Stage-C results, final claim. Show supported research claim separately from unsupported production and universal-quality claims.}
\caption{The paper claim is supported by frozen artifacts, while production and universal data-quality claims remain outside scope.}
\label{fig:evidence}
\end{figure}

\section{Discussion}
The result supports a bounded claim: in the studied code-domain setting, the
curation-stage framework produced a natural-budget curated arm that improved
held-out NLL and EvalPlus macro pass rate while using fewer packed training
tokens than raw fine-tuning.  However, identical LiveCodeBench pass@1 across
all three arms shows that this gain cannot yet be generalized to an independent
temporal code benchmark.  The math-domain validation fails under the
current selector, which is also evidence: the framework can expose that a
domain-specific Core/Metric/Policy configuration is not yet validated.  These
results do not imply that the framework measures intrinsic data quality, that
it will improve every future corpus, or that the risk detectors are
production-grade.

This boundary is important.  A raw corpus may already be sufficiently strong
after Stage 0 and Stage A, in which case forced removal would be harmful.  The
framework therefore allows retain-all outcomes.  Stage B is activated by a
budget constraint, not by an assumption that every corpus must be reduced.
Likewise, a corpus can be rejected or held for more validation if Stage-C
evidence is insufficient.

\subsection{Retain-All Is a Valid Outcome}
A common failure mode in automatic data curation is assuming that a framework
must always remove data.  That assumption is wrong for this setting.  If a
corpus passes Stage 0 and Stage A and fits the available training budget, the
correct framework output can be retain-all.  Stage B is a budget allocator, not
a moral judgment about discarded examples.  This design choice is important for
high-quality raw corpora, where unnecessary filtering can reduce coverage or
remove useful recurrence.

\subsection{Next Evidence Tier}
The next validation tier must test the same contract on corpora that mix
high-quality reference data with raw collected data.  The frozen protocol
requires equal-token fine-tuning arms for raw-random, Stage-A-random, curated,
and known-reference subsets; pre- and post-curation record, chunk, and token
counts; and external benchmark families matched to the domain.  For code, the
high-cost confirmatory target is SWE-bench-style issue resolution~\cite{swebench}
in addition to held-out code NLL and EvalPlus retention.  A powered multi-seed
LiveCodeBench run must also test whether the neutral 48-task pilot reflects low
power or a genuine lack of transfer.  For math, the target
benchmarks are GSM8K~\cite{gsm8k} and MATH~\cite{math}.  These experiments are
not claimed as completed evidence in this paper; they define the boundary
between the current code-domain result and a future multi-domain claim.

\section{Limitations and Threats to Validity}
\textbf{Core metric validity.}  The audits verify scoped operational behavior,
but they do not prove that every metric fully captures its intended construct.
Coverage remains limited by observable metadata, and Selection Value Evidence
is not a semantic quality oracle.

\textbf{Production readiness.}  The current Stage-0 detectors are project
benchmarks and fixtures, not production safety, legal, or license-compliance
certification.

\textbf{Domain scope.}  The strongest current training evidence is code-domain
evidence.  Math-domain evidence currently remains abstain after selector v3,
so the current selector must not be presented as a validated all-domain method.
Instruction, medical, legal, multilingual, and other specialized corpora
require their own frozen Stage-C validation before a training-use claim is
allowed.

\textbf{Feedback risk.}  Stage-C outcomes must not be used to tune Stage-B on
the same frozen holdouts.  The framework explicitly treats Utility as
validation-only to reduce this risk.

\section{Reproducibility}
The artifact set records commands, configs, source scripts, output artifacts,
documentation artifacts, and local hardware policy.  The current paper-evidence
package has a canonical lightweight rebuild path with seven scripts:
the Code evidence report, final evidence table, domain-composition audit,
coverage/domain-mix audit, Stage-B policy-contract audit, paper release gate,
and paper-claim consistency audit.  This path
rebuilds the claim-facing evidence package from existing Stage-C and guardrail
reports; it is deliberately not a full raw-data acquisition, benchmark
generation, GPU training, or production release runner.

\begin{table}[t]
\caption{Canonical Paper-Evidence Rebuild Path}
\label{tab:canonicalpath}
\centering
\footnotesize
\begin{tabular}{cl}
\toprule
Step & Script \\
\midrule
1 & \texttt{213\_build\_final\_paper\_evidence\_table.py} \\
2 & \texttt{219\_build\_domain\_composition\_audit.py} \\
3 & \texttt{220\_build\_coverage\_domain\_mix\_audit.py} \\
4 & \texttt{221\_build\_stage\_b\_policy\_contract\_audit.py} \\
5 & \texttt{218\_build\_paper\_claim\_consistency\_audit.py} \\
\bottomrule
\end{tabular}
\end{table}

The local default GPU policy records \texttt{CUDA\_VISIBLE\_DEVICES=1},
corresponding to the NVIDIA GeForce RTX 3070 Ti.  The RTX 4060 Ti is
approval-gated for local interactive-use reasons.  Production deployment
remains a separate blocked tier.

\section{Conclusion}
This paper presents a curation-stage framework for language-model training data
selection.  The framework separates Core responsibilities, measurable proxies,
and policy roles; enforces a staged boundary between hard gates, budgeted
selection, and downstream validation; and keeps Utility out of the selector
objective.  In the current evidence package, code is a positive natural-budget
validation case on NLL and EvalPlus, a LiveCodeBench pilot is neutral, and math
is an abstain case after a repaired selector still fails to beat raw training.
The claim is intentionally limited but useful:
the work operationalizes the data-curation stage for LM training, while
explicitly rejecting stronger claims about universal data-quality measurement,
all-domain improvement, or production-ready filtering.

\begin{thebibliography}{00}
\bibitem{ccnet}
G. Wenzek, M.-A. Lachaux, A. Conneau, V. Chaudhary, F. Guzman,
A. Joulin, and E. Grave, ``CCNet: Extracting high quality monolingual datasets
from web crawl data,'' in \emph{Proc. LREC}, 2020.

\bibitem{dedup}
K. Lee, D. Ippolito, A. Nystrom, C. Zhang, D. Eck, C. Callison-Burch,
and N. Carlini, ``Deduplicating training data makes language models better,''
in \emph{Proc. ACL}, 2022.

\bibitem{refinedweb}
G. Penedo, Q. Malartic, D. Hesslow, R. Cojocaru, A. Cappelli,
H. Alobeidli, B. Pannier, E. Almazrouei, and J. Launay, ``The RefinedWeb
dataset for Falcon LLM: Outperforming curated corpora with web data only,''
in \emph{Advances in Neural Information Processing Systems}, 2023.

\bibitem{dolma}
L. Soldaini et al., ``Dolma: An open corpus of three trillion tokens for
language model pretraining research,'' 2024.

\bibitem{finewebedu}
G. Penedo et al., ``The FineWeb datasets: Decanting the web for the finest text
data at scale,'' 2024.

\bibitem{c4}
C. Raffel et al., ``Exploring the limits of transfer learning with a unified
text-to-text transformer,'' \emph{Journal of Machine Learning Research}, 2020.

\bibitem{pile}
L. Gao et al., ``The Pile: An 800GB dataset of diverse text for language
modeling,'' 2020.

\bibitem{datacomp}
S. Y. Gadre et al., ``DataComp: In search of the next generation of multimodal
datasets,'' in \emph{Advances in Neural Information Processing Systems}, 2023.

\bibitem{swebench}
C. E. Jimenez, J. Yang, A. Wettig, S. Yao, K. Pei, O. Press, and
K. R. Narasimhan, ``SWE-bench: Can language models resolve real-world GitHub
issues?,'' in \emph{International Conference on Learning Representations}, 2024.

\bibitem{gsm8k}
K. Cobbe et al., ``Training verifiers to solve math word problems,'' 2021.

\bibitem{math}
D. Hendrycks et al., ``Measuring mathematical problem solving with the MATH
dataset,'' in \emph{Advances in Neural Information Processing Systems}, 2021.
\end{thebibliography}

\end{document}
